% Section 3. Budget 1.25 pages.
% Source: draft/bie-enclosure/s-reduction.tex (clean version).
\section{Reduction to the center cell}\label{sec:reduction}

Fix $\beta$. The center cell and the $j$th left and right periodic cells
($j\ge0$) are
\[
  \mathcal C_0=(x_-,x_+)\times I_d,\quad
  \mathcal C_{j,+}=\mathcal C_{0,+}+jL_+e_x,\quad
  \mathcal C_{j,-}=\mathcal C_{0,-}-jL_-e_x,
\]
with $\mathcal C_{0,+}=(x_+,x_++L_+)\times I_d$ and
$\mathcal C_{0,-}=(x_--L_-,x_-)\times I_d$; $B_\pm$ denotes the open half
strip formed by the cells on side $\pm$. The rods are $\Omega_0\subset\mathcal C_0$
(possibly empty) and $\Omega_{j,\pm}=\Omega_\pm\pm jL_\pm e_x$. The artificial
interfaces are $\Gamma_\pm=\{x_\pm\}\times I_d$ with normal $\nu_\pm=\pm e_x$,
pointing into $B_\pm$, and the walls $\Gamma_{j,\pm}=\{x_\pm\pm(j+1)L_\pm\}\times
I_d$ between consecutive cells. All normal derivatives on $\Gamma_\pm$ are taken
along $\nu_\pm$, from whichever side the function lives on.

For $O\in\{\mathcal C_0,B_-,B_+\}$ let $H^1_\beta(O)$ carry the quasi-periodic
condition in $y$, let $H^{1/2}_\beta(\Gamma_\pm)$ be its trace space and
$H^{-1/2}_\beta(\Gamma_\pm)$ the dual. For $u$ with $\Delta u\in L^2(O)$ and
quasi-periodic $\partial_yu$, the Cauchy data map
$\mathscr T_\pm u=(u|_{\Gamma_\pm},\partial_{\nu_\pm}u|_{\Gamma_\pm})$ is
defined weakly by Green's formula. The local solution spaces are
\begin{align*}
  \mathcal U_0(k)&=\{u\in H^1_\beta(\mathcal C_0):-\Delta u=k^2\rho u
  \text{ in }\mathcal C_0,\ \partial_yu\ \beta\text{-quasi-periodic}\},\\
  \mathcal V_\pm(k)&=\{v\in H^1_\beta(B_\pm):-\Delta v=k^2\rho_\pm v
  \text{ in }B_\pm,\ \partial_yv\ \beta\text{-quasi-periodic}\},
\end{align*}
with the equations in the distributional sense (so the rod transmission
conditions are included) and no condition on $\Gamma_\pm$. Functions in
$\mathcal V_\pm(k)$ have finite energy along the half strip. The
\emph{Cauchy data space} of the half guide is
$\mathcal G_\pm(k):=\mathscr T_\pm\mathcal V_\pm(k)$.

\begin{lemma}[unique continuation from Cauchy data]\label{lem:uc}
If $v\in\mathcal V_\pm(k)$ and $\mathscr T_\pm v=0$, then $v=0$.
\end{lemma}

\begin{proposition}[center reduction]\label{prop:reduction}
For fixed $\beta$ and $k>0$ let
$\mathcal U_0^{\mathcal G}(k)=\{u_0\in\mathcal U_0(k):\mathscr T_\pm u_0\in
\mathcal G_\pm(k)\}$. The restriction $u\mapsto u|_{\mathcal C_0}$ is a linear
bijection from $\ker(A(\beta)-k^2)$ onto $\mathcal U_0^{\mathcal G}(k)$.
\end{proposition}

The proofs are in \cref{sm:proof-reduction}: extending $v$ by zero across
$\Gamma_\pm$ through the homogeneous background strip gives an analytic
solution that vanishes on an open set; conversely, a center solution with
matching Cauchy data is glued to the unique half-guide solutions without a
distributional source on $\Gamma_\pm$.

\Cref{prop:reduction} holds for every $k>0$ and needs neither the gap nor
unique solvability of the half-guide Dirichlet problem. When the Dirichlet
problem is uniquely solvable, $\mathcal G_\pm(k)$ is the graph of the
Dirichlet-to-Neumann (DtN) operator and the reduction is the transparent
boundary condition of \cite[Thm.~4.5]{Fliss2013} (with the opposite normal).
At the countably many exceptional values where it is not
\cite[Thm.~4.1]{Fliss2013}, the DtN form fails while \cref{prop:reduction}
remains valid; Robin-to-Robin conditions avoid these values in another way
\cite{FlissKlindworthSchmidt2015}. In the gap, finite-energy half-guide
solutions decay cell by cell; Joly, Li, and Fliss represent them by a
propagation operator of spectral radius less than one \cite{JolyLiFliss2006},
whose discrete counterpart is the propagation matrix $\mathsf P_\pm$ of
\cref{sec:discretization}.
